\scp{Vowels}{12-05}
The reflexes of vowels in penultimate syllables in Banda are mostly straightforward.
The vowel *i is retained unchanged in nearly all instances,
as is *a in penultimate syllables.
(See further below for the reflexes of *a in final syllables.)
Similarly, *u is usually retained unchanged,
but penultimate *u > \it{i} when the following syllable contains
\it{i} (/{\gap}Ci) *duRi > \it{riri} ‘bone’,
*ma-udəhi > **mudi > \it{miri-} ‘back’,
and *kulit > \it{kilit} ‘skin’.

Antepenultimate vowels also assimilate either fully or partially
to the quality of the penultimate vowel: *p\tbr{a}niki > \it{f\tbr{i}niki}
‘bat’, *b\tbr{i}tuqən > \it{f\tbr{u}tuon} ‘star’,
*b\tbr{a}nua > \it{f\tbr{u}nua} ‘village’,
and *m\tbr{a}-tuqah > \it{m\tbr{o}tuo-, m\tbr{o}tuen} ‘old’.

In non-final syllables *ə > \it{e}.
Two words have *ə > \it{u} before \it{u} in the following syllable
(/{\gap}Cu) but there are also two examples of *ə > \it{e} in the same environment.
 Examples are given in \qf{12-12}.

\noindent{\wg\st{6pt}\tblr{>{\hspace{-\tabcolsep}}l<{\hspace{-\tabcolsep}}
>{\hspace{-\tabcolsep}\enspace}
lll}
\lsptoprule
%\tbx{12-12}	&	PMP	&	Banda	&	Banda gloss	\\	\midrule
%\endfirsthead
\tbx{12-12}	&	PMP	&	Banda	&	Banda gloss	\\	\midrule
%\endhead
	&	*haR\tbr{ə}zan	&	\tbr{e}ren	&	ladder	\\
	&	*ma-q\tbr{ə}ti	&	m\tbr{e}ti	&	ebb	\\
	&	*p\tbr{ə}ñu	&	\tbr{e}nu	&	tortoise	\\
	&	*qal\tbr{ə}jaw	&	l\tbr{e}a	&	sun	\\
	&	*qap\tbr{ə}ju	&	\tbr{e}lu-n	&	gall	\\
	&	*saŋ\tbr{ə}laR	&	ndes\tbr{e}nlar	&	fry	\\
	&	*t\tbr{ə}lu	&	t\tbr{e}lu	&	three	\\
	&	*qat\tbr{ə}luR	&	t\tbr{u}lur(u)	&	egg	\\
	&	*ma-p\tbr{ə}nuq	&	mb\tbr{u}nu	&	full	\\
\lspbottomrule
\end{tabular}\mbox{}\\}

Final *-uy > \it{u, o} as seen in *bab\tbr{uy} > \it{faf\tbr{u}} ‘pig’,
*hap\tbr{uy} > \it{a\tbr{u}} ‘fire’,
and *naŋ\tbr{uy} ‘swim’ > \it{nan\tbr{o}} ‘wade in neck deep water’.
We have only been able to identify one clear instance of *-iw:
*laR\tbr{iw} > ‹bĕ\it{ra}r\tbr{i}› ‘flee’.

The reflexes of *a
and *ə in final syllables,
as well as final vowel-glide sequences *-ay/*-aw are somewhat
complex with reflexes which partially overlap,
but do not fully merge.
To account for all the data we must keep final *ə,
*a, and *-ay/*-aw distinct.

In final syllables *ə > \it{o},
except after penultimate \it{e} in which case *ə > \it{e}.
PMP *ə only occurs in final closed syllables; /{\gap}C{\#}.
That is, there are no cases of *ə in final open syllables; /{\gap}{\#}.
Examples are given in \qf{12-13}.
(Transcriptions in angle brackets ‹ › are from \citet{vanEijbergen1864}
who uses italics to indicate stress.)

{\wg\begin{longtable}[l]{>{\hspace{-\tabcolsep}}
l<{\hspace{-\tabcolsep}}
>{\hspace{-\tabcolsep}\enspace}
llll}
\lsptoprule
\tbx{12-13}		&	env.	&	PMP	&	Banda	&	Banda gloss	\\	\midrule	\hhline{~}
\endfirsthead
\lsptoprule
\qf{12-13}		{\cwh}&	env.	{\cwh}&	PMP	{\cwh}&	Banda	{\cwh}&	Banda gloss	{\cwh}\\	\midrule
\endhead
&	u{\gap}{\#}		&	*bituq\tbr{ə}n	&	futu\tbr{o}n,	‹wĕ\it{toe}\tbr{a}n›&	star	\\
%\tbx{12-13}	&	u{\gap}{\#}		&	*bituq\tbr{ə}n	&	futu\tbr{o}n,					&	star	\\	\hhline{~}
	%&								&									&	‹wĕ\it{toe}\tbr{a}n›	&		\\
	&	u{\gap}{\#}		&	*put\tbr{ə}R		&	mbut\tbr{o}r					&	stir	\\
	&	aC{\gap}{\#}	&	*waR\tbr{ə}j		&	war\tbr{o}t						&	rope	\\
	&	aC{\gap}{\#}	&	*dal\tbr{ə}m		&	rar\tbr{o}-n(o)				&	inside	\\
	&	aC{\gap}{\#}	&	*tan\tbr{ə}m		&	ndan\tbr{o}						&	plant (v.)	\\
	&	aC{\gap}{\#}	&	*baRəq\tbr{ə}t	&	faf\tbr{o}t, faf\tbr{ó}t		&	heavy	\\	\hhline{~}
	%&	aC{\gap}{\#}	&	*baRəq\tbr{ə}t	&	faf\tbr{o}t,					&	heavy	\\	\hhline{~}
	%&								&									&	faf\tbr{ó}t						&		\\
	&	aC{\gap}{\#}	&	*qat\tbr{ə}p		&	at\tbr{o}p						&	roof thatch	\\
	&	aC{\gap}{\#}	&	*daR\tbr{ə}q		&	dar\tbr{o}						&	clay	\\
	&	aC{\gap}{\#}	&	*qaj\tbr{ə}ŋ		&	‹\it{a}l\tbr{o}n›			&	charcoal	\\
	&	oC{\gap}{\#}	&	*pus\tbr{ə}j		&	‹os\it{s}\tbr{o}loh› 	&	navel	\\
	&	iC{\gap}{\#}	&	*tiR\tbr{ə}m		&	‹ti\it{r}\tbr{on}›		&	oyster	\\
	&	eC{\gap}{\#}	&	*gət\tbr{ə}l		&	‹kètè\it{kè}t\tbr{e}›	&	itchy	\\
	&	eC{\gap}{\#}	&	*ma-qit\tbr{ə}m	&	mete{\tl}met\tbr{e}n,	mate{\tl}mat\tbr{e}n&	black	\\	\hhline{~}
	&	eC{\gap}{\#}	&	**mant\tbr{ə}R	&	m(b)end\tbr{e}r, 				&	cuscus	\\	\hhline{~}
	%&	eC{\gap}{\#}	&	*ma-qit\tbr{ə}m	&	mete{\tl}met\tbr{e}n,	&	black	\\	\hhline{~}
	%&								&									&	mate{\tl}mat\tbr{e}n	&		\\
	%&	eC{\gap}{\#}	&	**mant\tbr{ə}R	&	mend\tbr{e}r, 				&	cuscus	\\	\hhline{~}
	%&								&									&	mbend\tbr{e}r 				&		\\
\lspbottomrule
\end{longtable}}

In final closed syllables PMP *a = \it{a},
except after *z in which case *a > \it{e}.
This includes when the coda of the final syllable is a productive affix.
Examples are given in \qf{12-14}.
Note that \citet{Collins1986} gives some inalienably possessed
words with a final hyphen indicating that they obligatorily take
a possessive suffix, but does not give the form with the suffix, e.g.
*qabaRa > \it{fara-} ‘shoulder’.

{\wg\begin{longtable}[l]{>{\hspace{-\tabcolsep}}
l<{\hspace{-\tabcolsep}}
>{\hspace{-\tabcolsep}\enspace}
lll}
\lsptoprule
\tbx{12-14}	&	PMP	&	Banda	&	Banda gloss	\\	\midrule	\hhline{~}
\endfirsthead	
\lsptoprule
\qf{12-14}	{\cwh}&	PMP	{\cwh}&	Banda	{\cwh}&	Banda gloss	{\cwh}\\	\midrule
\endhead						
	&	*am\tbr{a}	&	am\tbr{a}-ŋ	&	father	\\	
	&	*an\tbr{a}k	&	an\tbr{a}ko	&	child	\\	
	&	*at\tbr{a}s	&	at\tbr{a}t	&	tall, high	\\	
	&	*bab\tbr{a}q	&	fof\tbr{á}n	&	below	\\	
	&	*bu\tbr{a}q	&	fu\tbr{a}-n	&	fruit, heart	\\	
	&	*bul\tbr{a}n	&	ful\tbr{a}n, wul\tbr{a}n	&	moon	\\	
	&	*daR\tbr{a}t	&	ndar\tbr{a}-k	&	go aground	\\	
	&	*əp\tbr{a}t	&	\tbr{a}t	&	four	\\	
	&	*habaR\tbr{a}t	&	far\tbr{a}t	&	west	\\	
	&	*hik\tbr{a}n	&	ik\tbr{a}n	&	fish	\\	
	&	*in\tbr{a}	&	in\tbr{a}-ŋ	&	mother	\\	
	&	*kak\tbr{a}k	&	kak\tbr{a}-n	&	older sibling	\\	
	&	**kamb\tbr{a}t	&	amb\tbr{a}t	&	wound	\\	
	&	**kəR\tbr{a}(ŋ)	&	ker\tbr{a}n	&	turtle	\\	
	&	*lay\tbr{a}R	&	le\tbr{a}r	&	sail	\\	
	&	*lu\tbr{a}q	&	lu\tbr{a}ko	&	vomit	\\	
	&	*mas\tbr{a}k	&	mosos\tbr{a}k	&	cook	\\	
	&	*mat\tbr{a}	&	mat\tbr{a}-n	&	eye	\\	
	&	*ñaR\tbr{a}	&	nar\tbr{a}-n	&	mother's sibling	\\	
	&	*ñaw\tbr{a}	&	‹njo\it{w\tbr{a}}na›	&	breathe	\\	
	&	*ŋaj\tbr{a}n	&	nal\tbr{a}n	&	name	\\	
	&	*pan\tbr{a}s	&	wan\tbr{a}t	&	hot	\\	
	&	*qalim\tbr{a}	&	lim\tbr{a}-n	&	hand	\\	
	&	*saŋ\tbr{a}	&	san\tbr{a-}, sen\tbr{a}-no	&	branch	\\	
	&	*saŋəl\tbr{a}R	&	ndesenl\tbr{a}r	&	fry	\\	
	&	*suq\tbr{a}n	&	su\tbr{a}n	&	digging stick	\\	
	&	*ti\tbr{a}n	&	ti\tbr{a}-n	&	belly	\\	
	&	*ud\tbr{a}ŋ	&	sugur\tbr{a}n	&	shrimp, prawn	\\	
	&	*uR\tbr{a}t	&	‹oer\tbr{a}t› /urat/	&	vein	\\	
	&	*wak\tbr{a}t	&	wak\tbr{a}t	&	root	\\	
	&	*haRəz\tbr{a}n	&	er\tbr{e}n	&	ladder	\\	
	&	*quz\tbr{a}n	&	ur\tbr{e}n	&	rain	\\	
\lspbottomrule
\end{longtable}}

In final open syllables,
including those that become open due to loss of a final consonant,
*a > \it{o {\tl} a}.
Some words have variant forms with final \it{o} and final \it{a}.
There are also two words which have *a > \it{e {\tl} a.} Examples
are given in \qf{12-15}.

\noindent{\wg\st{6pt}\tblr{>{\hspace{-\tabcolsep}}l<{\hspace{-\tabcolsep}}
>{\hspace{-\tabcolsep}\enspace}
lll}
\lsptoprule
\tbx{12-15}		&	PMP	&	Banda	&	Banda gloss	\\	\midrule
	&	*duha	&	ru\tbr{o}	&	two	\\
	&	*kita	&	kit\tbr{o}	&	\tsc{1pl.incl}	\\
	&	*kita	&	n-it\tbr{o}	&	see	\\
	&	*lima	&	lim\tbr{o}	&	five	\\
	&	*mamaq	&	mam\tbr{o}	&	chew	\\
	&	*panaq	&	mban\tbr{o}	&	shoot (bow)	\\
	&	*banua	&	funu\tbr{o}, funu\tbr{a}	&	village	\\
	&	*bayawak	&	mbuk\tbr{e}, mbuk\tbr{a}	&	monitor lizard	\\
	&	*daRaq	&	rar\tbr{o}-, rar\tbr{a}-	&	blood	\\
	&	*pija	&	il\tbr{o}, il\tbr{a}	&	how many?	\\
	&	*Rumaq	&	rum\tbr{o}, rum\tbr{a}	&	house	\\
	&	*ala[p/q]	&	n-al\tbr{a}	&	take	\\
	&	*ma-tuqah	&	motu\tbr{a}-, motu\tbr{e}n	&	old	\\
	&	*siwa	&	siw\tbr{a}	&	nine	\\
	&	*Raya	&	r\tbr{a}	&	big (likely /raa/)	\\
	&	*əsa	&	s\tbr{a}	&	one	\\
	&	*taq	&	t\tbr{a}	&	no	\\
\lspbottomrule
\end{tabular}\mbox{}\\}

At least some words also show a productive process of vowel alternation,
as shown in \qf{12-16},
exemplified with *Rumaq > \it{rumo {\tl} ruma} ‘house’.
The exact conditioning, if any, is unclear.
There appears to be a tendency for \it{o}-final forms to occur
as the citation form in lists,
but both \it{a} and \it{o} occur in phrase-final
and phrase-medial position.

\noindent{\wg\st{6pt}\begin{tabularx}
{\linewidth}{>{\hspace{-\tabcolsep}}l<{\hspace{-\tabcolsep}}
>{\hspace{-\tabcolsep}\enspace}l
>{\raggedright\arraybackslash}X}
\lsptoprule
%\tbx{12-16}	&	Banda	&	Banda gloss	\\	\midrule
%\endfirsthead
\tbx{12-16}	&	Banda	&	Banda gloss	\\	\midrule
%\endhead
	&	rum\tbr{o}	&	house	\\
	&	rum\tbr{o} sa	&	one house	\\
	&	rum\tbr{a} ra	&	big house	\\
	&	rum\tbr{a} aparok	&	uncle's house	\\
	&	rum\tbr{a} aparok ra	&	uncle's big house	\\
	&	kam ra-farak rum\tbr{a}	&	we (excl.) build a house	\\
	&	walin tari mutet rum\tbr{a}	&	(my) younger sibling has not yet occupied (his/her) house	\\
	&	kam mbelei ŋo wa rum\tbr{a}	&	we cannot leave (our) house	\\
	&	ɲaik ɲa rum\tbr{a}	&	depart from home	\\
\lspbottomrule
\end{tabularx}\mbox{}\\}

Vowel-glide sequences *-ay
and *-aw merge as \it{a} in nearly all cases.
No forms with a historical vowel-glide combination are known
to show variation between \it{a} and \it{o}.
Examples are given in \qf{12-17} and \qf{12-18}.
\citet{vanEijbergen1864} has one case of *-ay > ‹ai›; *matay
> ‹ma\it{tai}› ‘die’, ‹\it{ma}taité› ‘corpse’.
If the transcription in these words is accurate,
it may show that loss of final glides in Banda is a relatively
recent change.

\noindent{\wg\st{6pt}\tblr{>{\hspace{-\tabcolsep}}l<{\hspace{-\tabcolsep}}
>{\hspace{-\tabcolsep}\enspace}
lll}
\lsptoprule
%\tbx{12-17}	&	PMP	&	Banda	&	Banda gloss	\\	\midrule
%\endfirsthead
\tbx{12-17}		&	PMP	&	Banda	&	Banda gloss	\\	\midrule
%\endhead
&	*qaz\tbr{ay}	&	ar\tbr{a}-n	&	chin	\\
	&	*maRuqan\tbr{ay}	&	moran\tbr{a}, moron\tbr{a}	&	male	\\
	&	*mat\tbr{ay}	&	mat\tbr{a}-	&	dead	\\
	&	*nip\tbr{ay}	&	ni\tbr{a}, ɲi\tbr{a}	&	snake	\\
	&	*sak\tbr{ay}	&	ɲʤak\tbr{a}	&	ascend	\\
	&	*qat\tbr{ay}	&	at\tbr{a}-n	&	liver	\\
	&	*paj\tbr{ay}	&	al\tbr{a}	&	rice	\\
	&	*bəRs\tbr{ay}	&	mbes\tbr{e}, fes\tbr{e}	&	paddle (?**a > \it{e} /\it{e}C{\gap})	\\	\midrule
\end{tabular}\mbox{}\\}

\noindent{\wg\st{6pt}\tblr{>{\hspace{-\tabcolsep}}l<{\hspace{-\tabcolsep}}
>{\hspace{-\tabcolsep}\enspace}
lll}
\lsptoprule
%\tbx{12-17}	&	PMP	&	Banda	&	Banda gloss	\\	\midrule
%\endfirsthead
\tbx{12-18}		&	PMP	&	Banda	&	Banda gloss	\\	\midrule
%\endhead
&	**nak\tbr{aw}	&	mbo-nak\tbr{a}	&	steal	\\
	&	*balab\tbr{aw}	&	folaf\tbr{a}	&	mouse	\\
	&	*kas\tbr{aw}	&	kas\tbr{a}	&	rafter	\\
	&	*qaləj\tbr{aw}	&	le\tbr{a}	&	sun	\\
\lspbottomrule
\end{tabular}\mbox{}\\}

The reflexes of final vowels
and vowel-glide sequences are summarised in \qf{12-19} below.
To account for these reflexes,
we must propose *a > \it{o {\tl} a} occurred before loss of final glides.
This means that we cannot propose that *a merges with *-ay and *-aw.
The distinct reflexes of *ə also mean that we cannot propose
that *ə merged with *a or *-ay/*-aw.
Instead, *ə, *a, and *-ay/*-aw must be kept distinct to account
for the Banda data.
\mbox{}\\

\begin{tabular}{>{\hspace{-\tabcolsep}}l<{\hspace{-\tabcolsep}}ll@{ }l@{ }ll}
\tbx{12-19}	&	a.	&	*ə				&	>	&	\it{e}					&	/eC{\gap}C{\#}	\\
						&	b.	&	*ə				&	>	&	\it{o}					&	/{\gap}C{\#}	\\
						&	c.	&	*a				&	>	&	\it{a}					&	/{\gap}C{\#}	\\
						&	d.	&	*a				&	>	&	\it{o {\tl} a}	&	/{\gap}{\#}	\\
						&	e.	&	*-ay/*-aw	&	>	&	\it{a}					&	/{\gap}{\#}	\\
\end{tabular}
\mbox{}\\

Finally, in the literature two varieties of Banda are
identified named after the villages in which they
are spoken: Banda-Eli and Banda-Elat.
There is insufficient data for us to evaluate whether these
should be treated as two dialects of one language,
or two closely related languages.